Distributed systems that allow concurrent modifications to shared data face
the fundamental challenge of conflict resolution without central coordination.
The consistency, availability and partition tolerance (CAP) theorem, also known as
Brewer's Theorem \cite{CAP_Brewer}, dictates that no system can simultaneously
guarantee consistency, availability and partition tolerance \cite{CAP_main_aussage}.

Two predominant approaches have emerged:
Operational Transformation (OT) \cite{ellis1989concurrency}, originally developed
to support collaborative software, ensures consistency for concurrent editors.
Conflict-free Replicated Data Types (CRDTs) \cite{shapiro2011crdt} are data structures
replicated across multiple nodes that guarantee an application can update any replica
independently, concurrently, and without coordinating with other replicas.

Since guarentting all three properties simultaneously is impossible, every 
distributed system designer had to face the decision of which property to cut.
Many modern systems prioritize availability and partition tolerance over strong
consistency, accepting eventual consistency as a design goal \cite{vogels2009eventually}.

However, making a decision involves complex trade-offs in complexity, scalability
and architectural assumptions. 

Examples of this approach are found in collaborative editors 
like for an instance google docs or distributed databases.

Choosing the richt approach requires understanding the trade-offs between OT and
CRDTs in terms of complexity, scalability and architectural assumptions.

This paper analyzes and compares both approaches with respect to the following criteria.

The remainder of this paper is structered as follows: 
Section~\ref{sec:related}
reviews related work, Section~\ref{sec:methodology} describes the methodology,
Section~\ref{sec:concepts} presents the technical concepts, Section~\ref{sec:comparison}
provides a direct comparison, followed by a discussion in Section~\ref{sec:discussion}
and a conclusion in Section~\ref{sec:conclusion}.